% Chapter 1: motivation and the problem the thesis addresses.
\chapter{Introduction}\label{ch:intro}

During the July 2022 heatwave, air temperature at Aldmere's official
airport station peaked at 36.1\,°C. On the same afternoon, volunteer
sensors in the dense Old Harbour district recorded 39.4\,°C, and a
hospital two kilometres inland admitted three times its usual number of
patients with heat stroke \parencite{rautio2023harbour}. The official
record was not wrong; it simply measured a different place. Urban heat
varies over tens of metres, with street orientation, tree cover and
building height, while the networks that inform public warnings resolve
tens of kilometres \parencite{oke2017urban}.

Cheap connected sensors have made dense street-level networks
affordable, and many cities now run them
\parencite{chapman2017crowdsourcing,meier2017crowdsourcing}. Yet where
those sensors go is rarely a design decision. They are mounted where a
power outlet exists, where a school agreed to host one or where a
citizen-science group happens to be active. The result is a network that
is densest in affluent, well-organised districts and thinnest where
older residents live alone in top-floor flats, the population most at
risk \parencite{hsu2021disproportionate}.

\section{Problem statement}

This thesis addresses the question of \emph{where to put the next
sensor}. Given an existing network, a model of the temperature field and
a map of who is exposed, which site would most improve the information
used to protect people? Posed this way, sensor placement becomes an
instance of active learning: choosing the observation that is most
valuable under a model, then updating the model and choosing again
\parencite{krause2008near,settles2012active}.

Two features make the urban case harder than the textbook one. First,
value is not uniform: an error of one degree over an empty car park
matters less than the same error over a care home. Second, the model
must work before much local data exists, because a city deploying its
first network has nothing to learn from but other cities.

\section{Scope}

The thesis concerns near-surface air temperature during summer
afternoons and nights in mid-latitude European cities. Surface
temperature from satellites is used only as a covariate. Humidity and
thermal comfort indices are outside the main line of work, although the
field campaign records them for later use.
